\section{7.1 How to read this chapter}\label{how-to-read-this-chapter}

This chapter reports two instruments and the reader must not merge them.
\textbf{MGP-6 v4.3} (§7.10) is the successor release whose
construct-validity gate is enforced at build time; its sealed-hidden
result is the thesis's primary decision-level measurement, and its
development-split comparison is explicitly labelled development evidence
because that split selected the configuration. \textbf{MGP-6 v2}
(§7.2--§7.9) is the earlier pilot whose gate \emph{failed}; its numbers
are retained because the failure is itself a result and because it
motivated the v4.3 redesign, but no v2 row supports a capability claim.
\textbf{LongMemEval-M} (§7.11) is separate external validation, not a
substitute for either.

All values in this chapter are recomputed from released per-item
prediction files rather than transcribed from run reports; the
recomputation is archived in
\texttt{artifacts/analysis/mgp6\_v43\_full\_benchmark\_analysis.json}.

The construct-validity gate precedes every controller result. MGP-6 v2
is an author-designed pilot and its ledger label does not override the
newly measured shortcut failure. Policy scores below are diagnostic
observations on this collection, not estimates of general
memory-governance capability. Mapping-dependent external rows and every
LongMemEval row are excluded from primary comparisons.

\textbf{Shortcut probes fitted on train+development and evaluated on
test:}

{\def\LTcaptype{none} % do not increment counter
\begin{longtable}[]{@{}lrr@{}}
\toprule\noalign{}
Probe (no dialogue or memory state) & Accuracy & Macro-F1 \\
\midrule\noalign{}
\endhead
\bottomrule\noalign{}
\endlastfoot
Type only & 0.9667 & 0.9471 \\
Conflict type only & \textbf{1.0000} & \textbf{1.0000} \\
\end{longtable}
}

Conflict type deterministically identifies the action. This is the
chapter's primary result: the pilot cannot distinguish evidence
reasoning from scenario-family recognition. Cluster bootstrap quantifies
dependence inside the collection but cannot repair this construct
failure.

\section{7.2 MGP-6 v2 policy
diagnostics}\label{mgp-6-v2-policy-diagnostics}

Test split first (split-hygiene rule), full corpus second.

\textbf{Lexical reference policy.} Test split (n=90): action accuracy
\textbf{0.8889}, macro-F1 \textbf{0.9145}. Full corpus (n=180): 0.9222 /
0.9399. The released predictions were generated by deterministic lexical
predicates, not the prototype five-feature weighted scorecard. On test,
the reference makes 0/33 action false commits, with a one-sided 95\%
upper bound of 0.0868; it misses 10/45 required updates and
unnecessarily abstains on 10/57 non-ABSTAIN items.

\textbf{Structured LLM policy (qwen/qwen3.5-9b).} Test split: 0.9000 /
0.9040. Full corpus: 0.8889 / 0.8913. On test it makes 2/33 action false
commits and misses 7/45 updates. Its test accuracy is numerically above
the reference policy.

\textbf{Secondary full-corpus mechanism table (n = 180; includes train
and development):}

{\def\LTcaptype{none} % do not increment counter
\begin{longtable}[]{@{}
  >{\raggedright\arraybackslash}p{(\linewidth - 8\tabcolsep) * \real{0.2000}}
  >{\raggedright\arraybackslash}p{(\linewidth - 8\tabcolsep) * \real{0.2000}}
  >{\raggedright\arraybackslash}p{(\linewidth - 8\tabcolsep) * \real{0.2000}}
  >{\raggedright\arraybackslash}p{(\linewidth - 8\tabcolsep) * \real{0.2000}}
  >{\raggedright\arraybackslash}p{(\linewidth - 8\tabcolsep) * \real{0.2000}}@{}}
\toprule\noalign{}
\begin{minipage}[b]{\linewidth}\raggedright
System
\end{minipage} & \begin{minipage}[b]{\linewidth}\raggedright
Accuracy
\end{minipage} & \begin{minipage}[b]{\linewidth}\raggedright
Macro-F1
\end{minipage} & \begin{minipage}[b]{\linewidth}\raggedright
Invalid-commit
\end{minipage} & \begin{minipage}[b]{\linewidth}\raggedright
Eligibility
\end{minipage} \\
\midrule\noalign{}
\endhead
\bottomrule\noalign{}
\endlastfoot
Lexical reference policy & 0.9222 & 0.9399 & 0.0 & diagnostic \\
Structured LLM policy (qwen3.5-9b) & 0.8889 & 0.8913 & 0.0758 &
diagnostic \\
Zero-shot LLM decision & 0.700 & 0.6657 & 0.606 & diagnostic \\
B4 naive overwrite & 0.4222 & 0.2203 & 1.0 & diagnostic \\
B3 recency & 0.4222 & 0.2203 & 1.0 & diagnostic \\
B1 no-memory & 0.3667 & 0.1789 & --- & diagnostic \\
B2 flat retrieval & 0.1333 & 0.0784 & 1.0 & diagnostic \\
Few-shot (6-shot) decision & 0.7833 & 0.7692 & --- & incomplete: 174/180
predictions; 6 missing scored INVALID \\
\end{longtable}
}

The few-shot row is a \textbf{coverage-aware bound}, not a completed-run
comparison. Its 174 predictions were aligned to the 180 gold IDs and the
six missing IDs were scored INVALID. It is excluded from the
pre-specified significance family and from any claim that assumes
complete coverage.

\section{7.3 Ablations (n = 180; within-generator
diagnostics)}\label{ablations-n-180-within-generator-diagnostics}

{\def\LTcaptype{none} % do not increment counter
\begin{longtable}[]{@{}
  >{\raggedright\arraybackslash}p{(\linewidth - 6\tabcolsep) * \real{0.2500}}
  >{\raggedright\arraybackslash}p{(\linewidth - 6\tabcolsep) * \real{0.2500}}
  >{\raggedright\arraybackslash}p{(\linewidth - 6\tabcolsep) * \real{0.2500}}
  >{\raggedright\arraybackslash}p{(\linewidth - 6\tabcolsep) * \real{0.2500}}@{}}
\toprule\noalign{}
\begin{minipage}[b]{\linewidth}\raggedright
Configuration
\end{minipage} & \begin{minipage}[b]{\linewidth}\raggedright
Accuracy
\end{minipage} & \begin{minipage}[b]{\linewidth}\raggedright
Macro-F1
\end{minipage} & \begin{minipage}[b]{\linewidth}\raggedright
Distinct failure signature
\end{minipage} \\
\midrule\noalign{}
\endhead
\bottomrule\noalign{}
\endlastfoot
Full controller & 0.9222 & 0.9399 & --- \\
− abstention & 0.5556 & 0.4278 & invalid-commit 0.0 → \textbf{1.0} \\
− provenance & 0.700 & 0.7428 & invalid-commit → \textbf{0.606}
(identical to the zero-shot baseline's rate) \\
− typing & 0.7889 & 0.5680 & \textbf{KEEP F1 → 0.0} \\
− lineage & 0.9222 & 0.9399 & action labels unchanged; \textbf{lineage
integrity 0.8556 → 0.3689} \\
\end{longtable}
}

The lineage ablation is a useful within-generator mechanism check:
action labels remain identical while the instrument's lineage score
falls by more than half. It illustrates why decision and state metrics
should be reported separately, but author alignment and the pilot's
shortcut failure prevent treating the magnitude as an independent causal
estimate.

\section{7.4 Held-out clustered comparison and historical item
tests}\label{held-out-clustered-comparison-and-historical-item-tests}

On test, reference minus structured-LLM accuracy is \textbf{−0.0111}.
The paired template-family cluster interval \textbf{{[}−0.0938,
0.0625{]}} includes zero (48 test families); exact item-level McNemar
has b=4, c=5, p=1.0. The reference family-macro interval is {[}0.8125,
0.9688{]} and the structured-LLM interval is {[}0.8333, 0.9688{]}. These
data support neither superiority nor non-inferiority.

The table below preserves the earlier full-corpus item-level analysis
for audit history. It includes train and development items, assumes item
independence, and is not the primary inferential result.

{\def\LTcaptype{none} % do not increment counter
\begin{longtable}[]{@{}
  >{\raggedright\arraybackslash}p{(\linewidth - 12\tabcolsep) * \real{0.1429}}
  >{\raggedright\arraybackslash}p{(\linewidth - 12\tabcolsep) * \real{0.1429}}
  >{\raggedright\arraybackslash}p{(\linewidth - 12\tabcolsep) * \real{0.1429}}
  >{\raggedright\arraybackslash}p{(\linewidth - 12\tabcolsep) * \real{0.1429}}
  >{\raggedright\arraybackslash}p{(\linewidth - 12\tabcolsep) * \real{0.1429}}
  >{\raggedright\arraybackslash}p{(\linewidth - 12\tabcolsep) * \real{0.1429}}
  >{\raggedright\arraybackslash}p{(\linewidth - 12\tabcolsep) * \real{0.1429}}@{}}
\toprule\noalign{}
\begin{minipage}[b]{\linewidth}\raggedright
Pair
\end{minipage} & \begin{minipage}[b]{\linewidth}\raggedright
RD (A−B)
\end{minipage} & \begin{minipage}[b]{\linewidth}\raggedright
RD 95\% CI (Newcombe)
\end{minipage} & \begin{minipage}[b]{\linewidth}\raggedright
b / c
\end{minipage} & \begin{minipage}[b]{\linewidth}\raggedright
OR
\end{minipage} & \begin{minipage}[b]{\linewidth}\raggedright
Holm p
\end{minipage} & \begin{minipage}[b]{\linewidth}\raggedright
Decision
\end{minipage} \\
\midrule\noalign{}
\endhead
\bottomrule\noalign{}
\endlastfoot
Rule vs B4 naive overwrite & +0.5000 & {[}+0.4230, +0.5690{]} & 90 / 0 &
∞ & 4.85e-27 & reject H0 \\
Rule vs B2 flat retrieval & +0.7889 & {[}+0.7181, +0.8382{]} & 142 / 0 &
∞ & 1.44e-42 & reject H0 \\
Rule vs LLM controller & +0.0333 & {[}−0.0188, +0.0873{]} & 14 / 8 &
1.75 & \textbf{0.286} & \textbf{retain H0 (n.s.)} \\
LLM controller vs zero-shot & +0.1889 & {[}+0.1096, +0.2664{]} & 46 / 12
& 3.83 & 1.64e-05 & reject H0 \\
\end{longtable}
}

Interpretation, in the order that matters:

\begin{enumerate}
\def\labelenumi{\arabic{enumi}.}
\tightlist
\item
  \textbf{Reference vs structured LLM is a held-out null with the LLM
  numerically higher.} It is not evidence of equality, matching, or
  non-inferiority.
\item
  \textbf{Structured vs zero-shot is descriptive.} The historical
  full-corpus gap is large, but prompt structure and policy scaffolding
  are evaluated on a shortcut-leaking pilot; it is not a clean general
  scaffold-effect estimate.
\item
  \textbf{Reference vs B2/B4 is design-aligned.} The reference encodes
  the labeling manual while the baselines cannot express the required
  actions. These rows are mechanism checks, not capability proofs.
\end{enumerate}

\section{7.5 Archived external-system mapping sensitivity
(non-primary)}\label{archived-external-system-mapping-sensitivity-non-primary}

Under the favourable NOOP→ABSTAIN upper-bound mapping (NOOP→KEEP mapping
in parentheses):

{\def\LTcaptype{none} % do not increment counter
\begin{longtable}[]{@{}
  >{\raggedright\arraybackslash}p{(\linewidth - 6\tabcolsep) * \real{0.2500}}
  >{\raggedright\arraybackslash}p{(\linewidth - 6\tabcolsep) * \real{0.2500}}
  >{\raggedright\arraybackslash}p{(\linewidth - 6\tabcolsep) * \real{0.2500}}
  >{\raggedright\arraybackslash}p{(\linewidth - 6\tabcolsep) * \real{0.2500}}@{}}
\toprule\noalign{}
\begin{minipage}[b]{\linewidth}\raggedright
System
\end{minipage} & \begin{minipage}[b]{\linewidth}\raggedright
Accuracy
\end{minipage} & \begin{minipage}[b]{\linewidth}\raggedright
Macro-F1
\end{minipage} & \begin{minipage}[b]{\linewidth}\raggedright
Notes
\end{minipage} \\
\midrule\noalign{}
\endhead
\bottomrule\noalign{}
\endlastfoot
Letta (MemGPT lineage) \citep{packer2023memgpt} & 0.678 (0.489) & 0.488
(0.302) & \\
Mem0 2.0.11 \citep{chhikara2025mem0} & 0.622 (0.400) & 0.442 (0.249) &
invalid-commit 1.0 \\
Graphiti 0.29.2 \citep{rasmussen2025zep} & 0.383 (0.139) & 0.219 (0.104)
& \\
MemOS \citep{li2025memos} & 0.144 & 0.164 & \texttt{diagnostic\_only}:
supersession inexpressible through its adapter \\
\end{longtable}
}

These are harness- and mapping-specific transformations, not fair
primary capability measurements of the systems' native use cases. They
are excluded from the paper's main comparison. A confirmation study must
report native write/state/retrieval behavior first and show all action
mappings only as sensitivity analyses.

\section{\texorpdfstring{7.6 LongMemEval-M: diagnostic results
(\texttt{diagnostic\_only})}{7.6 LongMemEval-M: diagnostic results (diagnostic\_only)}}\label{longmemeval-m-diagnostic-results-diagnostic_only}

v7 stratified run (n = 180 each; judge google/gemma-4-31b-it):

{\def\LTcaptype{none} % do not increment counter
\begin{longtable}[]{@{}llll@{}}
\toprule\noalign{}
System & Correct & Accuracy & SSP \\
\midrule\noalign{}
\endhead
\bottomrule\noalign{}
\endlastfoot
B4 naive overwrite & 114 & \textbf{0.6333} & 11/30 \\
ASK controller & 110 & \textbf{0.6111} & 11/30 \\
B2 flat vector & 80 & 0.4444 & 6/30 \\
B3 recency & 43 & 0.2389 & 0/30 \\
B1 no-memory & 9 & 0.05 & --- \\
\end{longtable}
}

Two disclosures are mandatory. \textbf{First, the baseline ordering: on
this diagnostic slice, naive overwrite (63.3\%) outperformed the v7
controller (61.1\%).} Any claim of controller superiority on LongMemEval
is unsupported by current evidence. \textbf{Second, the v7 run is
provider-contaminated:} 316 HTTP-402 traces, 388 empty extraction
fallbacks cached as valid, at least 12 manual relaunches, and an
escalation tier that was a no-op (BIG\_MODEL == FAST\_MODEL). Controller
per-competency (v7, post-adjudication 110/180; pre-adjudication
108/180): knowledge-update 21/30, multi-session 15/30,
single-session-assistant 26/30, single-session-preference 11/30,
single-session-user 22/30, temporal-reasoning 15/30. Earlier
partial-range probes (v3: 5/10, 4/10; v4: 17/30; v6: 29/40) are
superseded diagnostics; v5 and the thinking-mode v8 replay are invalid;
both top-up runs are incomplete (§6.7).

\section{7.7 Why the failure-replay construction is
excluded}\label{why-the-failure-replay-construction-is-excluded}

The v8 replay re-executed exactly the 70 v7 controller failures under a
changed configuration and judge. It produced 25/70 machine-correct
items. A historical merge then retained all 110 v7 successes, replaced
only the replayed failures, and added four post-hoc Codex/GPT-5
overrides, yielding an archived 139/180 (0.7722) construction. That
number is classified \texttt{failure\_replay\_only}: the selection rule
makes judge and configuration identity perfectly confounded with prior
correctness, the fixes were designed on the same failures, the retained
cases inherited v7 contamination, and the four overrides were not
independent human adjudications. It is not a system result, an estimate
of final-1.1, or a confidence-bound endpoint.

The complementary Modal top-up makes the carry-forward error observable.
It targeted the 110 v7 successes and completed 84: \textbf{71 remained
correct and 13 regressed}. Combining these 84 observations with the 70
replay observations gives \textbf{96/154 raw observed correct}. The 26
unresolved items imply full-180 raw bounds of \textbf{96/180 = 53.3\%}
to \textbf{122/180 = 67.8\%}. If the four replay overrides are shown
separately as an adjudicated sensitivity, the corresponding bounds are
\textbf{55.6\%--70.0\%}. Thus this observed v8 attempt cannot reproduce
the archived carry-forward construction.

For SSP, the complementary evidence is equally important: the top-up
observes \textbf{8/29 correct}, with one unresolved item, so the maximum
is \textbf{9/30}. The v7 diagnostic 11/30 and the replay's 0/19 repairs
remain useful for causal investigation, but no SSP count can be carried
across configurations. The homogeneous final-1.1 full-180 run remains
\texttt{externally\_pending}: its smoke request failed before inference
with OpenRouter HTTP 403.

\section{7.8 Annotation agreement (inferential evidence about the
labels, not the
systems)}\label{annotation-agreement-inferential-evidence-about-the-labels-not-the-systems}

One independent annotator: raw agreement 0.9889 {[}0.9667, 1.0{]},
Cohen's κ 0.98146 {[}0.93997, 1.0{]}, one disagreement
(\texttt{MGP6V2-T3-TEST-016}). Four-LLM panel: Fleiss κ 0.9815, majority
vote 100\% vs gold, 88/90 unanimous. Caveats as in §5.5: n = 1 human +
model panel, not multi-human IAA.

\section{7.9 Hypothesis outcomes}\label{hypothesis-outcomes}

H1 is not identified by the shortcut-leaking pilot; H2 supports neither
superiority nor non-inferiority; H3 and H4 remain within-generator
mechanism demonstrations; H5 is unsupported. The independently generated
v3 confirmation protocol is the required gate for stronger conclusions.

\section{7.10 MGP-6 v4.3 final hidden result (one frozen
configuration)}\label{mgp-6-v4.3-final-hidden-result-one-frozen-configuration}

\subsection{7.10.1 Headline result and its
uncertainty}\label{headline-result-and-its-uncertainty}

The pre-registered ASK live controller attained \textbf{0.8481 accuracy}
and \textbf{0.8479 macro-F1} on all 540 sealed hidden items, coverage
1.000. The component-cluster bootstrap (5,000 resamples, seed 4301)
gives 95\% intervals of \textbf{{[}0.8037, 0.8889{]}} for accuracy and
\textbf{{[}0.8021, 0.8894{]}} for macro-F1.

Because the hidden split is exactly action-balanced (180 KEEP, 180
SUPERSEDE, 180 ABSTAIN), two agreement statistics are interpretable
without prior correction, and both are reported here for the first time:

{\def\LTcaptype{none} % do not increment counter
\begin{longtable}[]{@{}lrl@{}}
\toprule\noalign{}
Quantity & Value & 95\% interval \\
\midrule\noalign{}
\endhead
\bottomrule\noalign{}
\endlastfoot
Action accuracy & 0.8481 & {[}0.8037, 0.8889{]} component-cluster \\
--- item-independence reference & & {[}0.8154, 0.8760{]} Wilson \\
Macro-F1 & 0.8479 & {[}0.8021, 0.8894{]} component-cluster \\
Balanced accuracy & 0.8481 & --- \\
Cohen's κ vs gold & \textbf{0.7722} & --- \\
Matthews correlation (multiclass) & \textbf{0.7728} & --- \\
Non-abstention coverage & 0.6926 & --- \\
Selective accuracy & 0.8449 & --- \\
\end{longtable}
}

Two observations follow. First, the component-cluster interval is
\textbf{41\% wider} than the Wilson interval that item independence
would imply (width 0.0852 versus 0.0605). That difference is the cost of
the component design, and reporting only the Wilson width would
understate the uncertainty by 29\%. Second, the gap between 0.848
accuracy and 0.772 κ is entirely the \emph{structure} of the errors,
since chance accuracy is exactly 1/3 by construction.

Per-action precision / recall / F1 is ABSTAIN 0.8554 / 0.7889 / 0.8208,
KEEP 0.7947 / 0.8389 / 0.8162, and SUPERSEDE 0.8967 / 0.9167 / 0.9066,
each on support 180. Prediction marginals are 166 ABSTAIN, 190 KEEP, 184
SUPERSEDE.

\subsection{7.10.2 Confusion matrix and the error
decomposition}\label{confusion-matrix-and-the-error-decomposition}

{\def\LTcaptype{none} % do not increment counter
\begin{longtable}[]{@{}lrrr@{}}
\toprule\noalign{}
Gold ~Predicted & ABSTAIN & KEEP & SUPERSEDE \\
\midrule\noalign{}
\endhead
\bottomrule\noalign{}
\endlastfoot
ABSTAIN & \textbf{142} & 30 & 8 \\
KEEP & 18 & \textbf{151} & 11 \\
SUPERSEDE & 6 & 9 & \textbf{165} \\
\end{longtable}
}

The errors are strongly asymmetric, and for a write-governance system
the asymmetry --- not the total --- is the operational result. SUPERSEDE
is the best-recovered action (0.9167 recall): when the evidence
genuinely establishes a new value, the controller applies it. ABSTAIN is
the worst (0.7889 recall), and its 38 failures split 30 to KEEP and only
8 to SUPERSEDE.

That split is the thesis's central measurement point about error
taxonomy. An ABSTAIN→KEEP error leaves the stored state
\textbf{byte-identical} to what abstention would have produced; the
memory is not corrupted, and only the audit record that the system saw
ambiguous evidence and declined is lost. An ABSTAIN→SUPERSEDE error
\textbf{actually replaces an active record on insufficient evidence} ---
the precise failure the three-action space exists to prevent. Collapsing
both into a single ``false commit'' rate of 38/180 = 0.2111 therefore
overstates the harm by roughly fivefold relative to the state-corrupting
component.

{\def\LTcaptype{none} % do not increment counter
\begin{longtable}[]{@{}
  >{\raggedright\arraybackslash}p{(\linewidth - 6\tabcolsep) * \real{0.2143}}
  >{\raggedleft\arraybackslash}p{(\linewidth - 6\tabcolsep) * \real{0.2857}}
  >{\raggedleft\arraybackslash}p{(\linewidth - 6\tabcolsep) * \real{0.2857}}
  >{\raggedright\arraybackslash}p{(\linewidth - 6\tabcolsep) * \real{0.2143}}@{}}
\toprule\noalign{}
\begin{minipage}[b]{\linewidth}\raggedright
Governance error
\end{minipage} & \begin{minipage}[b]{\linewidth}\raggedleft
Count / denominator
\end{minipage} & \begin{minipage}[b]{\linewidth}\raggedleft
Rate
\end{minipage} & \begin{minipage}[b]{\linewidth}\raggedright
Wilson 95\% CI
\end{minipage} \\
\midrule\noalign{}
\endhead
\bottomrule\noalign{}
\endlastfoot
Action false commit (gold ABSTAIN acted on) & 38/180 & 0.2111 &
{[}0.1578, 0.2764{]} \\
--- of which ABSTAIN→KEEP (state preserved) & 30/180 & 0.1667 &
{[}0.1193, 0.2280{]} \\
--- of which \textbf{ABSTAIN→SUPERSEDE (state corrupted)} &
\textbf{8/180} & \textbf{0.0444} & \textbf{{[}0.0227, 0.0852{]}} \\
Missed update (gold SUPERSEDE not applied) & 15/180 & 0.0833 &
{[}0.0511, 0.1329{]} \\
Destructive overwrite (gold KEEP superseded) & 11/180 & 0.0611 &
{[}0.0345, 0.1061{]} \\
Unnecessary abstention (gold KEEP/SUPERSEDE abstained) & 24/360 & 0.0667
& {[}0.0452, 0.0973{]} \\
\end{longtable}
}

Unnecessary abstention at 0.0667 is low, which matters: the controller
is not purchasing its conservatism by refusing to act. But selective
accuracy (0.8449) is \emph{below} overall accuracy (0.8481), so the
items it does commit on are answered slightly less accurately than
average. A well-calibrated selective predictor shows the opposite
ordering. \textbf{The abstention mechanism is therefore doing governance
work --- avoiding unjustified writes --- without doing selective-risk
work}, and this thesis does not conflate the two. This is the clearest
single defect in the frozen configuration and it is disclosed rather
than omitted.

\subsection{7.10.3 Component decomposition: which part does the
work}\label{component-decomposition-which-part-does-the-work}

The controller composes a public-evidence rule policy with a live model
policy under a conservative guard that overrides a model ABSTAIN only
when the rule policy finds decisive public evidence. Each component is
separately scorable on development data.

{\def\LTcaptype{none} % do not increment counter
\begin{longtable}[]{@{}lrr@{}}
\toprule\noalign{}
& Public dev (n=180) & Hidden (n=540) \\
\midrule\noalign{}
\endhead
\bottomrule\noalign{}
\endlastfoot
Model action alone & 0.7500 & not scorable† \\
Rule policy alone & 0.7500 & not scorable† \\
\textbf{Composed (reported)} & \textbf{0.8444} & \textbf{0.8481} \\
Items where the guard changed the action & 27 & 36 \\
--- repaired a model error & 20 & --- \\
--- broke a correct model action & 3 & --- \\
--- error-neutral & 4 & --- \\
\end{longtable}
}

† Component-wise hidden scoring would require a second pass over sealed
gold beyond the single authorised scoring event.

On the hidden split the guard fired 36 times and always out of ABSTAIN
(30 → KEEP, 6 → SUPERSEDE). This \textbf{bounds} its contribution: even
under the most favourable accounting, the guard accounts for at most 6.7
accuracy points of the 84.8 observed. The rules that fired were
no-decisive-public-evidence (255), direct-documented-slot-change (120),
standing-value-confirmed-or-proposal- retracted (90),
unresolved-scope-authority-or-timing (57), and
temporary-or-scoped-exception (18). The result is a genuine composition,
not a rule engine with a model attached nor a model with cosmetic rules.

\subsection{7.10.4 Cost sensitivity across governance
regimes}\label{cost-sensitivity-across-governance-regimes}

{\def\LTcaptype{none} % do not increment counter
\begin{longtable}[]{@{}lrrr@{}}
\toprule\noalign{}
System & Balanced & Commit-averse & Update-responsive \\
\midrule\noalign{}
\endhead
\bottomrule\noalign{}
\endlastfoot
Frozen controller, \textbf{hidden} (n=540) & 0.1519 & 0.4204 & 0.2963 \\
Frozen controller, dev & 0.1556 & 0.4167 & 0.3000 \\
Successor controller, dev & 0.0500 & 0.1500 & 0.1333 \\
Successor + composed guard, dev & 0.0389 & 0.0889 & 0.0944 \\
B4 naive overwrite, dev & 0.6611 & 3.2333 & 1.6222 \\
Rule policy alone, dev & 0.6667 & 1.2444 & 1.6667 \\
\end{longtable}
}

The commit-averse loss is the largest for the frozen controller because
that regime charges most heavily for exactly its dominant error class.
Critically, \textbf{the ordering among ASK controllers is stable across
all three regimes} --- unlike the v2 reference-versus-LLM comparison
(§7.4), where the preferred system \emph{flips} with the cost matrix.
This is a stronger result than the v2 chapter could produce and it is
stated as such.

\subsection{7.10.5 Computation time and token
cost}\label{computation-time-and-token-cost}

All runs used concurrency 1, temperature 0, and a 128-token generation
cap, so these are per-item serial costs and not a throughput claim.

{\def\LTcaptype{none} % do not increment counter
\begin{longtable}[]{@{}
  >{\raggedright\arraybackslash}p{(\linewidth - 14\tabcolsep) * \real{0.0968}}
  >{\raggedleft\arraybackslash}p{(\linewidth - 14\tabcolsep) * \real{0.1290}}
  >{\raggedleft\arraybackslash}p{(\linewidth - 14\tabcolsep) * \real{0.1290}}
  >{\raggedleft\arraybackslash}p{(\linewidth - 14\tabcolsep) * \real{0.1290}}
  >{\raggedleft\arraybackslash}p{(\linewidth - 14\tabcolsep) * \real{0.1290}}
  >{\raggedleft\arraybackslash}p{(\linewidth - 14\tabcolsep) * \real{0.1290}}
  >{\raggedleft\arraybackslash}p{(\linewidth - 14\tabcolsep) * \real{0.1290}}
  >{\raggedleft\arraybackslash}p{(\linewidth - 14\tabcolsep) * \real{0.1290}}@{}}
\toprule\noalign{}
\begin{minipage}[b]{\linewidth}\raggedright
Run
\end{minipage} & \begin{minipage}[b]{\linewidth}\raggedleft
n
\end{minipage} & \begin{minipage}[b]{\linewidth}\raggedleft
mean (ms)
\end{minipage} & \begin{minipage}[b]{\linewidth}\raggedleft
median
\end{minipage} & \begin{minipage}[b]{\linewidth}\raggedleft
p95
\end{minipage} & \begin{minipage}[b]{\linewidth}\raggedleft
max
\end{minipage} & \begin{minipage}[b]{\linewidth}\raggedleft
prompt tok/item
\end{minipage} & \begin{minipage}[b]{\linewidth}\raggedleft
completion tok/item
\end{minipage} \\
\midrule\noalign{}
\endhead
\bottomrule\noalign{}
\endlastfoot
Frozen controller, hidden & 540 & 1,673 & 1,278 & 2,931 & 12,626 & 393 &
7.4 \\
Frozen controller, dev & 180 & 1,400 & 959 & 3,669 & 10,260 & 396 &
7.5 \\
Successor controller, dev & 180 & 1,288 & 952 & 3,388 & 7,593 & 506 &
7.8 \\
Zero-shot decision, dev & 180 & 1,425 & 1,039 & 3,122 & 9,886 & 549 &
8.0 \\
B4 naive overwrite, dev & 180 & 1,647 & 1,139 & 4,116 & 23,586 & 578 &
7.9 \\
B2 flat retrieval, dev & 180 & 1,581 & 1,060 & 4,582 & 8,478 & 597 &
7.9 \\
B3 recency, dev & 180 & 1,774 & 1,071 & 5,515 & 11,404 & 233 & 7.9 \\
Rule policy (no model call), dev & 180 & 0.24 & 0.13 & 0.64 & 3.6 & 0 &
0 \\
\end{longtable}
}

The write decision is cheap. Hidden-split totals are 212,467 prompt
tokens, 4,000 completion tokens, and 903.5 s of summed model latency
across 540 items. The latency distribution is right-skewed with a
\emph{thin} tail rather than a fat one: median-to-p95 is a factor of 2.3
while p95-to-max is 4.3, so the worst cases are route-level stalls, not
hard items. These are provider-reported wall times through a shared
OpenRouter route and are \textbf{not} compute measurements; the
deterministic rule policy, which issues no network call, completes the
same 180 decisions at a 0.13 ms median.

Token cost separates the systems more informatively than latency.
Evidence-indexed routing spends 393 prompt tokens per hidden decision
against 578 for naive overwrite, 597 for flat retrieval, and 549 for
zero-shot --- all of which score at chance. The successor configuration
costs 28\% more than the frozen one (506 vs 396 prompt tokens on dev)
and buys 10.6 accuracy points.

\textbf{Monetary cost is unavailable, not estimated.} OpenRouter
telemetry carried no price or cost field on any run in this thesis.
Given the frozen model identifier a list-price multiplication would be
arithmetically trivial and epistemically empty: routing, provider
selection, and caching all affect the billed amount and none of that
appears in the logs. Token counts are recorded because they were
measured.

Run integrity on the hidden split: 0 in-client retries, 0 invalid or
unparseable responses, 0 duplicate predictions, and one transport-level
disconnect after 432 predictions recovered by the immutable-manifest
resume under identical controls. One caveat on the manifest --- its
207.96 s \texttt{wall\_clock\_seconds} field records only the resumed
invocation that produced the final 108 predictions, so it is
\textbf{not} the end-to-end duration of the 540-item run. The 903.5 s
summed model latency is the defensible aggregate.

\subsection{7.10.6 Development-split system comparison (n = 180,
balanced)}\label{development-split-system-comparison-n-180-balanced}

Because only one hidden configuration was authorised, every system
contrast is development evidence on the split used for selection. Chance
accuracy is 0.3333.

{\def\LTcaptype{none} % do not increment counter
\begin{longtable}[]{@{}
  >{\raggedright\arraybackslash}p{(\linewidth - 14\tabcolsep) * \real{0.1000}}
  >{\raggedleft\arraybackslash}p{(\linewidth - 14\tabcolsep) * \real{0.1333}}
  >{\raggedright\arraybackslash}p{(\linewidth - 14\tabcolsep) * \real{0.1000}}
  >{\raggedleft\arraybackslash}p{(\linewidth - 14\tabcolsep) * \real{0.1333}}
  >{\raggedleft\arraybackslash}p{(\linewidth - 14\tabcolsep) * \real{0.1333}}
  >{\raggedleft\arraybackslash}p{(\linewidth - 14\tabcolsep) * \real{0.1333}}
  >{\raggedleft\arraybackslash}p{(\linewidth - 14\tabcolsep) * \real{0.1333}}
  >{\raggedleft\arraybackslash}p{(\linewidth - 14\tabcolsep) * \real{0.1333}}@{}}
\toprule\noalign{}
\begin{minipage}[b]{\linewidth}\raggedright
System
\end{minipage} & \begin{minipage}[b]{\linewidth}\raggedleft
Acc.
\end{minipage} & \begin{minipage}[b]{\linewidth}\raggedright
95\% CI
\end{minipage} & \begin{minipage}[b]{\linewidth}\raggedleft
Macro-F1
\end{minipage} & \begin{minipage}[b]{\linewidth}\raggedleft
MCC
\end{minipage} & \begin{minipage}[b]{\linewidth}\raggedleft
FC
\end{minipage} & \begin{minipage}[b]{\linewidth}\raggedleft
MU
\end{minipage} & \begin{minipage}[b]{\linewidth}\raggedleft
DO
\end{minipage} \\
\midrule\noalign{}
\endhead
\bottomrule\noalign{}
\endlastfoot
Successor controller (guard-composed) & \textbf{0.9611} & {[}0.933,
0.989{]} & 0.9607 & 0.943 & 0.000 & 0.017 & 0.033 \\
Successor controller (as run) & 0.9500 & {[}0.917, 0.978{]} & 0.9494 &
0.926 & 0.000 & 0.033 & 0.067 \\
Frozen controller (pre-registered) & 0.8444 & {[}0.789, 0.894{]} &
0.8458 & 0.771 & 0.233 & 0.133 & 0.033 \\
B1 no memory & 0.3667 & {[}0.294, 0.439{]} & 0.3118 & 0.065 & 0.800 &
0.217 & 0.733 \\
B2 flat retrieval & 0.3556 & {[}0.283, 0.428{]} & 0.2440 & 0.065 & 0.900
& 0.067 & 0.917 \\
B3 recency & 0.3278 & {[}0.261, 0.400{]} & 0.2175 & −0.015 & 0.933 &
0.117 & 0.900 \\
B4 naive overwrite & 0.3389 & {[}0.272, 0.411{]} & 0.1891 & 0.023 &
1.000 & 0.017 & 0.967 \\
Zero-shot LLM decision & 0.3389 & {[}0.272, 0.411{]} & 0.1788 & 0.046 &
1.000 & 0.000 & 0.983 \\
Rule policy alone (no model) & 0.3333 & {[}0.267, 0.406{]} & 0.2434 &
0.000 & 0.217 & 0.783 & 0.217 \\
− abstention & 0.3333 & {[}0.267, 0.406{]} & 0.2465 & 0.000 & 1.000 &
0.233 & 0.767 \\
− provenance & 0.3389 & {[}0.272, 0.411{]} & 0.1887 & 0.029 & 0.983 &
0.017 & 0.967 \\
− typing & 0.3444 & {[}0.278, 0.417{]} & 0.1899 & 0.092 & 0.983 & 0.000
& 0.983 \\
− lineage & 0.3500 & {[}0.283, 0.422{]} & 0.2020 & 0.088 & 0.983 & 0.000
& 0.967 \\
Letta (adapter-mapped) & 0.2944 & {[}0.228, 0.361{]} & 0.2424 & −0.069 &
0.417 & 0.733 & 0.383 \\
Mem0 (adapter-mapped) & 0.3333 & {[}0.267, 0.406{]} & 0.1851 & 0.000 &
0.033 & 1.000 & 0.000 \\
Graphiti (adapter-mapped) & 0.3056 & {[}0.239, 0.372{]} & 0.2112 &
−0.064 & 0.217 & 0.867 & 0.117 \\
\end{longtable}
}

FC = action false commit on gold-ABSTAIN; MU = missed update on
gold-SUPERSEDE; DO = destructive overwrite of a gold-KEEP record.

The most instructive fact here is negative. \textbf{Every} baseline,
ablation, zero-shot policy, standalone rule policy, and deployed
external system lands between 0.2944 and 0.3667 --- at or below chance
--- with macro-F1 in 0.179--0.312 and MCC in −0.069 to 0.092. They are
close to uninformative on this task, and they fail in two
\emph{opposite} ways. Commit-happy systems (B2, B4, zero-shot, and the
no-abstention/no-typing/no-lineage ablations) supersede almost
everything: destructive-overwrite rates 0.767--0.983, false-commit rates
0.900--1.000. Abstain-happy systems (rule policy alone, Mem0, Graphiti)
do the reverse: Mem0 abstains on 117/120 items requiring a decision and
misses all 60 required updates. Both score near chance while behaving in
opposite ways --- which is precisely why this thesis reports the error
decomposition rather than accuracy alone.

Paired exact McNemar tests, Holm-corrected over the 15-test family:

{\def\LTcaptype{none} % do not increment counter
\begin{longtable}[]{@{}
  >{\raggedright\arraybackslash}p{(\linewidth - 12\tabcolsep) * \real{0.1200}}
  >{\raggedright\arraybackslash}p{(\linewidth - 12\tabcolsep) * \real{0.1200}}
  >{\raggedleft\arraybackslash}p{(\linewidth - 12\tabcolsep) * \real{0.1600}}
  >{\raggedleft\arraybackslash}p{(\linewidth - 12\tabcolsep) * \real{0.1600}}
  >{\raggedleft\arraybackslash}p{(\linewidth - 12\tabcolsep) * \real{0.1600}}
  >{\raggedleft\arraybackslash}p{(\linewidth - 12\tabcolsep) * \real{0.1600}}
  >{\raggedright\arraybackslash}p{(\linewidth - 12\tabcolsep) * \real{0.1200}}@{}}
\toprule\noalign{}
\begin{minipage}[b]{\linewidth}\raggedright
A
\end{minipage} & \begin{minipage}[b]{\linewidth}\raggedright
B
\end{minipage} & \begin{minipage}[b]{\linewidth}\raggedleft
b
\end{minipage} & \begin{minipage}[b]{\linewidth}\raggedleft
c
\end{minipage} & \begin{minipage}[b]{\linewidth}\raggedleft
RD
\end{minipage} & \begin{minipage}[b]{\linewidth}\raggedleft
Holm p
\end{minipage} & \begin{minipage}[b]{\linewidth}\raggedright
Decision
\end{minipage} \\
\midrule\noalign{}
\endhead
\bottomrule\noalign{}
\endlastfoot
Frozen controller & Letta & 106 & 7 & +0.5500 & 1.2e−22 & reject H0 \\
Frozen controller & − abstention & 98 & 6 & +0.5111 & 2.2e−21 & reject
H0 \\
Frozen controller & B1 no memory & 90 & 4 & +0.4778 & 4.2e−21 & reject
H0 \\
Frozen controller & B4 naive overwrite & 98 & 7 & +0.5056 & 1.4e−20 &
reject H0 \\
Frozen controller & Rule policy alone & 106 & 14 & +0.5111 & 5.8e−18 &
reject H0 \\
Frozen controller & Mem0 & 106 & 14 & +0.5111 & 5.8e−18 & reject H0 \\
Successor (guard-comp.) & Frozen controller & 28 & 7 & +0.1167 & 1.5e−03
& reject H0 \\
Successor (as run) & Frozen controller & 28 & 9 & +0.1056 & 5.1e−03 &
reject H0 \\
Successor (guard-comp.) & Successor (as run) & 2 & 0 & +0.0111 & 0.500 &
\textbf{retain H0} \\
\end{longtable}
}

(Remaining contrasts against B2, B3, zero-shot, − provenance, − typing,
− lineage, and Graphiti all reject at Holm p \textless{} 1.6e−19; full
table in the released analysis artifact.)

The guard-composition repair therefore did \textbf{not} measurably
improve the successor system; it restored the intended composition
semantics without damaging accuracy, and this thesis claims no more than
that.

\subsection{7.10.7 What this result is
not}\label{what-this-result-is-not}

This is one measurement of one configuration on a controlled synthetic
diagnostic after public-development model selection. It is \textbf{not}
a comparison against the development-only baselines or ablations under
sealed conditions, not a paired hidden test, and not evidence of
real-world memory-governance performance or deployment safety. The
successor configuration's 0.9611 development accuracy has not been
tested on any sealed split and must not be read as an improvement over
the 0.8481 hidden result: they are different splits under different
exposure regimes. LongMemEval (§7.11) remains a separate
external-validation track.

\subsection{7.11 LongMemEval-M V12: the homogeneous full-slice
run}\label{longmemeval-m-v12-the-homogeneous-full-slice-run}

The V12 run (freeze \texttt{final-2.0}) evaluates all 180 stratified
items with one answer model (\texttt{qwen/qwen3.5-9b}), one judge
(\texttt{google/gemma-4-31b-it}), and zero judge-model fallbacks. It is
the eligible downstream result on this track.

{\def\LTcaptype{none} % do not increment counter
\begin{longtable}[]{@{}
  >{\raggedright\arraybackslash}p{(\linewidth - 10\tabcolsep) * \real{0.1364}}
  >{\raggedleft\arraybackslash}p{(\linewidth - 10\tabcolsep) * \real{0.1818}}
  >{\raggedright\arraybackslash}p{(\linewidth - 10\tabcolsep) * \real{0.1364}}
  >{\raggedleft\arraybackslash}p{(\linewidth - 10\tabcolsep) * \real{0.1818}}
  >{\raggedleft\arraybackslash}p{(\linewidth - 10\tabcolsep) * \real{0.1818}}
  >{\raggedleft\arraybackslash}p{(\linewidth - 10\tabcolsep) * \real{0.1818}}@{}}
\toprule\noalign{}
\begin{minipage}[b]{\linewidth}\raggedright
System
\end{minipage} & \begin{minipage}[b]{\linewidth}\raggedleft
Correct/n
\end{minipage} & \begin{minipage}[b]{\linewidth}\raggedright
Accuracy (95\% CI)
\end{minipage} & \begin{minipage}[b]{\linewidth}\raggedleft
Prompt tok/item
\end{minipage} & \begin{minipage}[b]{\linewidth}\raggedleft
Latency/item (s)
\end{minipage} & \begin{minipage}[b]{\linewidth}\raggedleft
vs.~row above
\end{minipage} \\
\midrule\noalign{}
\endhead
\bottomrule\noalign{}
\endlastfoot
B1 no memory & 9/180 & 0.0500 {[}0.022, 0.083{]} & 95 & 1.6 & --- \\
B3 recency & 43/180 & 0.2389 {[}0.183, 0.300{]} & 768 & 3.1 & 4.7e−10 \\
B2 flat vector & 80/180 & 0.4444 {[}0.372, 0.517{]} & 3,629 & 2.5 &
3.7e−06 \\
B4 naive overwrite & 114/180 & 0.6333 {[}0.561, 0.700{]} & 14,227 &
966.0 & 6.8e−06 \\
\textbf{ASK governed writes} & \textbf{137/180} & \textbf{0.7611}
{[}0.700, 0.822{]} & 22,625 & 809.2 & \textbf{2.9e−04} \\
\end{longtable}
}

Every adjacent step in the monotone ordering separates under an exact
paired McNemar test, Holm-corrected over the six tests reported. The
controller-over- naive-overwrite contrast is b = 31, c = 8, RD = +0.128,
Holm p = 2.9e−04.

Per competency (30 items each):

{\def\LTcaptype{none} % do not increment counter
\begin{longtable}[]{@{}
  >{\raggedright\arraybackslash}p{(\linewidth - 10\tabcolsep) * \real{0.1304}}
  >{\raggedleft\arraybackslash}p{(\linewidth - 10\tabcolsep) * \real{0.1739}}
  >{\raggedleft\arraybackslash}p{(\linewidth - 10\tabcolsep) * \real{0.1739}}
  >{\raggedleft\arraybackslash}p{(\linewidth - 10\tabcolsep) * \real{0.1739}}
  >{\raggedleft\arraybackslash}p{(\linewidth - 10\tabcolsep) * \real{0.1739}}
  >{\raggedleft\arraybackslash}p{(\linewidth - 10\tabcolsep) * \real{0.1739}}@{}}
\toprule\noalign{}
\begin{minipage}[b]{\linewidth}\raggedright
Competency
\end{minipage} & \begin{minipage}[b]{\linewidth}\raggedleft
B1
\end{minipage} & \begin{minipage}[b]{\linewidth}\raggedleft
B3
\end{minipage} & \begin{minipage}[b]{\linewidth}\raggedleft
B2
\end{minipage} & \begin{minipage}[b]{\linewidth}\raggedleft
B4
\end{minipage} & \begin{minipage}[b]{\linewidth}\raggedleft
ASK
\end{minipage} \\
\midrule\noalign{}
\endhead
\bottomrule\noalign{}
\endlastfoot
knowledge update & 0.133 & 0.667 & 0.567 & 0.700 & \textbf{0.967} \\
multi-session & 0.033 & 0.200 & 0.133 & 0.567 & \textbf{0.633} \\
single-session assistant & 0.000 & 0.000 & 0.900 & 0.867 &
\textbf{0.967} \\
single-session preference & 0.000 & 0.000 & 0.200 & 0.367 &
\textbf{0.533} \\
single-session user & 0.067 & 0.400 & 0.767 & 0.767 & \textbf{0.833} \\
temporal reasoning & 0.067 & 0.167 & 0.100 & 0.533 & \textbf{0.633} \\
\end{longtable}
}

The gain over naive overwrite is \textbf{not uniform}, and its shape is
the strongest available evidence that the improvement comes from the
write decision rather than from generic pipeline differences. It
concentrates in knowledge update (+8 items) and single-session
preference (+5) --- exactly where a stale value must be replaced or a
durable preference separated from a transient one. On multi-session the
gain is +2 items and on temporal reasoning +3, both within noise at this
slice size. A better retriever would not preferentially help
knowledge-update items.

\textbf{Write-decision behaviour.} The controller's action log records
12,654 write decisions across 3,547 sessions: 11,478 KEEP (90.7\%), 779
ABSTAIN (6.2\%), and 397 SUPERSEDE (3.1\%) over candidate types
episodic\_fact (9,857), transient\_pref (2,416), and stable\_pref (381).
This distribution is what the design predicts on natural dialogue, where
most turns carry no conflicting evidence, and it differs sharply from
the balanced MGP-6 v4.3 diagnostic --- which is the reason both
instruments are retained. All 397 supersessions recorded both a
predecessor and a successor identifier, and none of the 779 abstentions
produced a resulting record, so the lineage and no-op invariants held on
all 12,654 decisions.

\subsection{Pipeline anatomy: where the compute actually
goes}\label{pipeline-anatomy-where-the-compute-actually-goes}

The released pipeline traces (\texttt{extraction\_trace.jsonl},
\texttt{relation\_trace.jsonl}, \texttt{answers.jsonl}) permit a
stage-level decomposition rather than a single per-item cost figure, and
the decomposition overturns the natural reading of the accuracy table.

{\def\LTcaptype{none} % do not increment counter
\begin{longtable}[]{@{}
  >{\raggedright\arraybackslash}p{(\linewidth - 10\tabcolsep) * \real{0.1304}}
  >{\raggedleft\arraybackslash}p{(\linewidth - 10\tabcolsep) * \real{0.1739}}
  >{\raggedleft\arraybackslash}p{(\linewidth - 10\tabcolsep) * \real{0.1739}}
  >{\raggedleft\arraybackslash}p{(\linewidth - 10\tabcolsep) * \real{0.1739}}
  >{\raggedleft\arraybackslash}p{(\linewidth - 10\tabcolsep) * \real{0.1739}}
  >{\raggedleft\arraybackslash}p{(\linewidth - 10\tabcolsep) * \real{0.1739}}@{}}
\toprule\noalign{}
\begin{minipage}[b]{\linewidth}\raggedright
Stage
\end{minipage} & \begin{minipage}[b]{\linewidth}\raggedleft
Model calls
\end{minipage} & \begin{minipage}[b]{\linewidth}\raggedleft
Prompt tok.
\end{minipage} & \begin{minipage}[b]{\linewidth}\raggedleft
Completion tok.
\end{minipage} & \begin{minipage}[b]{\linewidth}\raggedleft
Share of prompt
\end{minipage} & \begin{minipage}[b]{\linewidth}\raggedleft
Model latency
\end{minipage} \\
\midrule\noalign{}
\endhead
\bottomrule\noalign{}
\endlastfoot
Extraction & 11,847 & 41,806,968 & 1,806,685 & \textbf{73.6\%} & 19.98
h \\
Relation classification & 19,157 & 10,936,290 & 1,226,921 &
\textbf{19.2\%} & 18.93 h \\
Answer generation & 180 & 4,072,477 & 41,418 & \textbf{7.2\%} & see
note \\
\textbf{Total} & \textbf{31,184} & \textbf{56,815,735} &
\textbf{3,075,024} & --- & 38.91 h \\
Per item & 173.2 & 315,643 & 17,083 & --- & 778 s \\
\end{longtable}
}

Answering the 180 questions takes 180 model calls and 4.07 M prompt
tokens. Reaching the point where those questions \emph{can} be answered
takes 31,004 further calls and 52.7 M prompt tokens. The write-time
stages are \textbf{92.8\% of prompt tokens and 99.4\% of model calls};
the answer stage --- the only stage an end-to-end memory benchmark
scores --- is 7.2\% of the budget.

This reframes the cost comparison above. The governed controller is not
an expensive answerer; it is a cheap answerer sitting on an expensive
ingestion pipeline, and the same holds for naive overwrite. Reporting
only answer-stage tokens, as the retrieval-only baselines invite,
measures the 7\% and ignores the 93\%.

\textbf{A correction to how per-item latency must be read.} The
\texttt{latency\_ms} field on each answer record is \textbf{cumulative
end-to-end} time for that item, not answer-generation time. Summing it
gives 40.46 h, while summed extraction and relation model latency is
38.91 h; the 1.55 h difference is the implied generation-only cost ---
about \textbf{31 s per item}, not the 809 s the field suggests. Per-item
pipeline times quoted in this section use the cumulative field and must
be read as end-to-end pipeline time. This thesis flags the point
explicitly because the naive reading --- that one answer call takes
thirteen minutes --- is wrong by more than an order of magnitude.

\textbf{Extraction is mostly rejection.} Of 11,491 processed turns,
\textbf{5,656 (49.2\%)} return \texttt{no\_extractable\_content}. The
extractor's most common output is \emph{nothing}, which is correct
behaviour on conversational filler and the first line of defence against
memory bloat. The remaining turns yield 11,771 candidate memories (1.02
per turn overall), typed as 9,144 episodic facts, 2,261 transient
preferences, and 366 stable preferences. Mean extraction confidence is
0.934 (minimum 0.75); 17.4\% of candidates carry a hedge marker and
1.3\% a negation marker. Parsing was robust: 370 of 11,491 turns hit a
first-parse error, 9 required the repair path, and there were
\textbf{zero model fallbacks and zero retries beyond the first attempt}.

\textbf{Relation routing is deliberately conservative --- and
auditable.} The stage is two-part: stage one labels the candidate
against a retrieved prior as \texttt{support}, \texttt{contradict}, or
\texttt{independent}; stage two, invoked on 9,530 of 9,539
classifications, resolves same-subject and temporal ordering. Routing is
then deterministic and every decision records its reason:

{\def\LTcaptype{none} % do not increment counter
\begin{longtable}[]{@{}llr@{}}
\toprule\noalign{}
Stage-1 label & Routed relation & n \\
\midrule\noalign{}
\endhead
\bottomrule\noalign{}
\endlastfoot
independent & unclear & 8,281 \\
support & support & 880 \\
\textbf{contradict} & \textbf{supersede} & \textbf{221} \\
contradict & contradict & 147 \\
contradict & unclear & 1 \\
parse fallback & unclear & 9 \\
\end{longtable}
}

All 8,281 \texttt{independent} labels route to \texttt{unclear} --- an
independent candidate is never permitted to displace a prior. Of 369
\texttt{contradict} labels, only the 221 that stage two confirms as
\textbf{same-subject and newer} route to \texttt{supersede}; 113 older
and 33 same-time contradictions are held as \texttt{contradict} without
a write. So \textbf{221 of 9,539 classifications (2.3\%)} may propose a
replacement. This is the mechanism behind the 397 supersessions (3.1\%)
in the action log and, ultimately, behind the controller's low
destructive-overwrite rate --- the conservatism is architectural, not
emergent.

\textbf{Caching did not help, and the trace explains why.} The
extraction cache holds 11,477 unique keys against 11,491 traced turns,
the relation cache 9,523 against 9,539 classifications --- under 0.3\%
deduplication in both. Each turn is keyed by its own content and each
candidate by its pairing with a specific prior, so in a single pass over
distinct sessions there is essentially nothing to reuse. The caches are
resume insurance, not a throughput optimisation, and any deployment
estimate assuming cache amortisation across a first pass is unsupported
by this evidence.

\textbf{Downstream cost in context.} Per-item end-to-end pipeline time
is 809 s for the governed controller and 966 s for naive overwrite,
against 1.6--3.1 s for the three retrieval-only baselines --- roughly
320× against flat-vector retrieval, buying 31.7 accuracy points.
Answer-stage prompt tokens show the same ordering (22,625 per item
versus 14,227 and 3,629), but as the stage decomposition above shows,
that comparison understates the true gap by more than an order of
magnitude, because the retrieval-only baselines have no ingestion
pipeline to account for. Whether the exchange is worthwhile is a
deployment question this thesis does not adjudicate; the full accounting
is reported so it is not hidden behind the accuracy ordering.

\textbf{Judge accounting.} Every item was scored by the single frozen
judge: 40 items resolved by deterministic guards (normalised match, no
LLM call), 140 routed to the LLM judge, 2 scoring fallbacks logged, 0
model fallbacks. Nine items received a documented semantic-equivalence
adjudication correcting judge scoring errors, each with a written
rationale and a preserved pre-adjudication backup.

\textbf{Two open provenance items} do not affect any reported number:
the homogeneity gate flags a stale code/protocol hash for the evaluation
launcher, and the manifest was not written to a clean terminal status
after a resume. Per-system integrity checks (answer model, judge
routing, fallback counts) pass with empty error and warning lists.
